\chapter{Cerebral Aneurysm Clipping and Coiling}

\section*{Introduction \& Clinical Indications}
Cerebral aneurysm treatment excludes a weak arterial outpouching from circulation by microsurgical clipping or endovascular coiling, flow diversion, or other embolization. Treatment is considered after subarachnoid hemorrhage, or for selected unruptured aneurysms whose rupture risk exceeds the treatment risk. Decisions depend on aneurysm size, location, morphology, growth, symptoms, age, family history, prior hemorrhage, comorbidity, and patient preference.

Clipping requires craniotomy and placement of a clip across the aneurysm neck. Endovascular treatment uses a catheter introduced through an artery to place coils or another device. The choice is multidisciplinary; neither method is universally best.

\section*{Relevant Surgical Anatomy}
The circle of Willis and its branches supply the brain and cranial nerves. Aneurysm location determines exposure, perforator risk, parent-vessel preservation, and access route. The subarachnoid cisterns, optic nerves, chiasm, brainstem, and venous structures may be at risk during clipping. Endovascular planning evaluates the parent artery, branch vessels, neck width, access tortuosity, and collateral circulation.

\section*{Preoperative Workup \& Preparation}
CTA, MRA, and catheter angiography define the aneurysm. Ruptured cases require assessment for hydrocephalus, vasospasm, rebleeding, seizures, electrolyte abnormalities, and systemic instability. Consent includes stroke, hemorrhage, vessel occlusion, cranial-nerve injury, seizures, cognitive change, recurrence, need for retreatment, contrast reaction, and death.

\section*{Surgical Technique \& Procedure}
Clipping involves craniotomy, microsurgical dissection, temporary or permanent vascular control, and clip placement confirmed by angiography or Doppler. Coiling advances a microcatheter into the aneurysm and deploys coils, sometimes with balloon, stent, or flow-diverter assistance. Antiplatelet therapy may be required for some devices. Final imaging assesses aneurysm occlusion and parent-vessel patency.

\section*{Complications \& Risks}
Risks include intraoperative rupture, thromboembolism, ischemic stroke, hemorrhage, vasospasm, hydrocephalus, seizures, infection, cranial-nerve injury, brain swelling, contrast or access complications, coil migration, incomplete occlusion, recurrence, and delayed rebleeding.

\section*{Postoperative Care \& Recovery}
Neurocritical care monitors consciousness, focal deficits, blood pressure, sodium, fluid balance, seizures, and vasospasm. Ruptured aneurysm patients require surveillance for delayed cerebral ischemia and hydrocephalus. Rehabilitation may address cognition, speech, swallowing, mobility, and mood. Follow-up angiography is tailored to treatment, aneurysm anatomy, and residual filling.

\section*{Outcomes \& Prognosis}
Prognosis is dominated by hemorrhage severity and early brain injury in ruptured cases. Complete durable occlusion is not guaranteed, particularly for wide-neck or complex aneurysms. Long-term imaging and management of hypertension, smoking, and other vascular risks are important.

\section*{References}
\begin{enumerate}
  \item American Heart Association/American Stroke Association. Guideline for the Management of Patients With Aneurysmal Subarachnoid Hemorrhage.
  \item European Stroke Organisation. Guidelines on Management of Unruptured Intracranial Aneurysms.
  \item Congress of Neurological Surgeons. Guidelines for Treatment of Cerebral Aneurysms.
\end{enumerate}

\clearpage
